\chapter{CÁLCULOS ELÉCTRICOS}
\section{Conductor y longitud eléctrica}
Se adopta conductor AAAC 6201-T81 de sección nominal 50 mm$^2$, sección real de cálculo 49.5 mm$^2$, diámetro 9 mm y masa 135 kg/km. La longitud por fase se obtiene de la catenaria EDS de cada vano; no se reemplaza por la proyección horizontal.
\begin{align}
L_{\mathrm{XY}}&=\sum_{i=1}^{110}L_i=6\,696.369\ \mathrm{m},\\
\ell&=\sum_{i=1}^{110}L_{c,i}=6\,758.120\ \mathrm{m}.
\end{align}
\section{Resistencia y reactancia}
\begin{align}
R_{75}&=R_{20}[1+\alpha_R(75-20)]\frac{\ell}{1000}\\
&=0.671[1+0.0036(55)](6.758120)\notag\\
&=\boxed{5.432569\ \Omega/\mathrm{fase}}.\notag
\end{align}
Con separaciones horizontales de 1.2, 1.2 y 2.4 m:
\begin{align}
D_m&=(D_{ab}D_{bc}D_{ca})^{1/3}=1.511905\ \mathrm{m},\\
\mathrm{GMR}&=0.7788\frac{0.009}{2}=0.0035046\ \mathrm{m},\\
X&=2\pi f(2\times10^{-4})\ln\!\left(\frac{D_m}{\mathrm{GMR}}\right)\frac{\ell}{1000}\\
&=\boxed{3.091467\ \Omega/\mathrm{fase}}.\notag
\end{align}
\section{Modelo de transformación}
Cada transformador se representa por impedancia serie y consumo en vacío. Se emplean $S_n=250$ kVA, $u_k=0.06$ pu, $P_{cu,n}=3.1$ kW, $P_0=0.45$ kW y $Q_0=2.5$ kvar.
\begin{align}
r_T&=\frac{P_{cu,n}}{S_n}=0.0124\ \mathrm{pu},\\
x_T&=\sqrt{u_k^2-r_T^2}=0.058705\ \mathrm{pu},\\
Z_b&=\frac{V_n^2}{S_n}=400\ \Omega,\\
Z_T&=(r_T+jx_T)Z_b=4.960+j23.482\ \Omega.
\end{align}
Las pérdidas y reactancias son datos tipo del modelo. La elevación efectiva de fuente $t=1.05$ pu representa una condición equivalente de regulación, sin atribuirle una posición física de toma no documentada.
\section{Flujo de potencia con T1, línea y T2}
El modelo trifásico equilibrado se resuelve por fase mediante barrido atrás/adelante. La tensión de fuente se fija en $V_0=10\,000(1.05)/\sqrt{3}$ V y las cargas se representan como potencias PQ:
\begin{align}
S_1&=S_2=450+j2500\ \mathrm{VA},\\
S_3&=171\,486+j83\,054.460\ \mathrm{VA},\\
I_j^{(k)}&=\sum_{m=j}^{3}\left(\frac{S_m}{3V_m^{(k)}}\right)^{*},\\
\widehat V_{j+1}^{(k+1)}&=\widehat V_j^{(k+1)}-Z_jI_j^{(k)},\\
V_j^{(k+1)}&=0.7\widehat V_j^{(k+1)}+0.3V_j^{(k)}.
\end{align}
La iteración se detiene con diferencia de tensión inferior a $10^{-8}$ V. El libro de cálculo desarrolla 60 iteraciones visibles, con partes real e imaginaria, sin referencias circulares.
\begingroup\small
\begin{longtable}{L{\dimexpr 0.666667\linewidth-2\tabcolsep\relax}L{\dimexpr 0.333333\linewidth-2\tabcolsep\relax}}
\toprule\rowcolor{modelblue!12}
\textbf{Magnitud} & \textbf{Resultado} \\ \midrule
\endfirsthead
\toprule\rowcolor{modelblue!12}
\textbf{Magnitud} & \textbf{Resultado} \\ \midrule
\endhead
MT de salida de T1 & 10.181 kV \\
MT de llegada a T2 & 10.059 kV \\
BT de llegada & 371.067 V \\
Corriente en la línea & 11.356330 A \\
Pérdida activa de línea & 2.101854 kW \\
Generación requerida & 183.326151 kW \\
Potencia aparente de fuente & 212.480501 kVA \\
\bottomrule\end{longtable}\endgroup

\begin{align}
\Delta P_{\mathrm{LP}}&=3|I_2|^2R=2.101854\ \mathrm{kW},\\
\Delta V_{\mathrm{LP}}&=\frac{|V_1|-|V_2|}{|V_1|}\,100=1.207337\%,\\
V_{\mathrm{BT}}&=\sqrt{3}|V_3|\frac{380}{10\,000}=371.066630\ \mathrm{V}.
\end{align}
La caída anterior corresponde sólo a la línea. No debe confundirse con la regulación total generador--T1--línea--T2. Con fuente nominal equivalente sin elevación, la BT calculada es 350.291 V; la diferencia justifica analizar la regulación del conjunto.
\section{Comparación de tensiones}
\begingroup\small
\begin{longtable}{L{\dimexpr 0.2\linewidth-2\tabcolsep\relax}L{\dimexpr 0.2\linewidth-2\tabcolsep\relax}L{\dimexpr 0.2\linewidth-2\tabcolsep\relax}L{\dimexpr 0.2\linewidth-2\tabcolsep\relax}L{\dimexpr 0.2\linewidth-2\tabcolsep\relax}}
\toprule\rowcolor{modelblue!12}
\textbf{Tensión nominal kV} & \textbf{Fuente pu} & \textbf{BT llegada V} & \textbf{Pérdida LP kW} & \textbf{Fuente kW} \\ \midrule
\endfirsthead
\toprule\rowcolor{modelblue!12}
\textbf{Tensión nominal kV} & \textbf{Fuente pu} & \textbf{BT llegada V} & \textbf{Pérdida LP kW} & \textbf{Fuente kW} \\ \midrule
\endhead
10.0 & 1.000 & 350.291 & 2.359 & 184.052 \\
10.0 & 1.025 & 360.709 & 2.224 & 183.672 \\
10.0 & 1.050 & 371.067 & 2.102 & 183.326 \\
20.0 & 1.000 & 354.320 & 0.576 & 182.172 \\
20.0 & 1.025 & 364.601 & 0.544 & 181.906 \\
20.0 & 1.050 & 374.832 & 0.515 & 181.663 \\
22.0 & 1.000 & 354.550 & 0.476 & 182.066 \\
22.0 & 1.025 & 364.823 & 0.449 & 181.806 \\
22.0 & 1.050 & 375.047 & 0.425 & 181.569 \\
22.9 & 1.000 & 354.634 & 0.439 & 182.027 \\
22.9 & 1.025 & 364.905 & 0.414 & 181.770 \\
22.9 & 1.050 & 375.126 & 0.392 & 181.534 \\
\bottomrule\end{longtable}\endgroup

El aumento de tensión reduce la corriente y las pérdidas de la línea. No elimina el déficit de potencia activa de una fuente aislada ni demuestra compatibilidad de equipos existentes. La alternativa 22.9 kV se conserva como contraste; el caso de cálculo de este expediente permanece identificado en 10 kV.
